\documentclass[12pt, a4paper]{article}

% ---- Standard Packages ----
\usepackage[utf8]{inputenc}
\usepackage[T1]{fontenc}
\usepackage[bahasa]{babel}
\usepackage{amsmath, amssymb, amsthm}
\usepackage{geometry}
\usepackage{enumitem}
\usepackage{booktabs}

\geometry{margin=1in}

\title{BANK SOAL MATEMATIKA KELAS 10\\ \large Latihan Mandiri Persiapan UAS}
\author{Ringkasan Latihan Materi}
\date{}

\begin{document}

\maketitle

% ==============================================================
\section{EKSPONEN (BILANGAN PANGKAT)}
% ==============================================================
Sederhanakan atau tentukan nilai dari bentuk-bentuk berikut:
\begin{enumerate}
    \item Hasil dari $2^5 \times 2^3 \times 2^{-4}$ adalah ...
    \item Bentuk sederhana dari $\frac{a^3 b^2}{a^{-1} b^4}$ adalah ...
    \item Nilai dari $(3^2)^3 \times 3^{-5}$ adalah ...
    \item Sederhanakan $\left( \frac{2x^2 y^{-3}}{4x^{-1} y^2} \right)^2$.
    \item Tentukan nilai $x$ yang memenuhi $3^{x+1} = 81$.
    \item Hitunglah nilai dari $125^{2/3} + 64^{1/3} - 81^{1/4}$.
    \item Jika $a=2$ dan $b=3$, tentukan nilai dari $\frac{a^2 b^3}{a^{-1} b}$.
    \item Ubahlah ke dalam bentuk pangkat positif: $\frac{7x^{-2} y^5}{z^{-3}}$.
    \item Sederhanakan $(2^3 \cdot 3^2)^2$.
    \item Tentukan nilai dari $\frac{10^6 \times 4^2}{25^3 \times 8^3}$.
\end{enumerate}

% ==============================================================
\section{BENTUK AKAR}
% ==============================================================
Sederhanakan atau rasionalkan bentuk-bentuk berikut:
\begin{enumerate}
    \item Sederhanakan $\sqrt{72} + \sqrt{50} - \sqrt{18}$.
    \item Hitunglah $3\sqrt{5} \times 2\sqrt{10}$.
    \item Hasil dari $(\sqrt{3} + \sqrt{2})(\sqrt{3} - \sqrt{2})$ adalah ...
    \item Rasionalkan penyebut dari $\frac{10}{\sqrt{5}}$.
    \item Rasionalkan penyebut dari $\frac{6}{3 - \sqrt{3}}$.
    \item Tentukan bentuk sederhana dari $\frac{\sqrt{3}}{\sqrt{5} + \sqrt{2}}$.
    \item Sederhanakan $\sqrt[3]{16} \times \sqrt[3]{4}$.
    \item Hitunglah $2\sqrt{27} - \sqrt{75} + 3\sqrt{12}$.
    \item Rasionalkan $\frac{4}{2\sqrt{3}}$.
    \item Tentukan nilai dari $\sqrt{48} : \sqrt{3}$.
\end{enumerate}

% ==============================================================
\section{LOGARITMA}
% ==============================================================
Gunakan sifat-sifat logaritma untuk menyelesaikan soal berikut:
\begin{enumerate}
    \item Nilai dari ${}^2\log 16 + {}^2\log 4 - {}^2\log 2$ adalah ...
    \item Jika ${}^3\log 2 = a$, tentukan nilai dari ${}^3\log 8$.
    \item Tentukan nilai dari ${}^5\log 100 - {}^5\log 4$.
    \item Hitunglah ${}^2\log 3 \cdot {}^3\log 5 \cdot {}^5\log 8$.
    \item Sederhanakan ${}^a\log \frac{1}{b} \cdot {}^b\log \frac{1}{c} \cdot {}^c\log \frac{1}{a}$.
    \item Tentukan nilai $x$ dari ${}^x\log 49 = 2$.
    \item Jika $\log 2 = 0,301$ dan $\log 3 = 0,477$, hitunglah $\log 6$.
    \item Tentukan nilai dari ${}^8\log 4$.
    \item Sederhanakan ${}^2\log 25 \cdot {}^5\log 3 \cdot {}^3\log 32$.
    \item Tentukan nilai $n$ jika ${}^2\log (n+1) = 4$.
\end{enumerate}

% ==============================================================
\section{BARISAN DAN DERET}
% ==============================================================
\subsection{Aritmatika}
\begin{enumerate}
    \item Tentukan suku ke-20 dari barisan $3, 7, 11, 15, \dots$
    \item Diketahui barisan aritmatika dengan $U_3 = 10$ dan $U_7 = 22$. Tentukan beda ($b$) dan suku pertama ($a$).
    \item Hitunglah jumlah 15 suku pertama dari deret $2 + 5 + 8 + 11 + \dots$
    \item Suatu barisan memiliki rumus $U_n = 5n - 3$. Tentukan suku ke-50.
    \item Dalam sebuah aula terdapat 10 baris kursi. Baris terdepan ada 20 kursi, baris berikutnya selalu bertambah 4 kursi. Berapa total kursi di aula tersebut?
\end{enumerate}

\subsection{Geometri}
\begin{enumerate}
    \item Tentukan suku ke-6 dari barisan $2, 6, 18, \dots$
    \item Diketahui barisan geometri dengan $U_1 = 5$ dan $r = 2$. Tentukan $U_8$.
    \item Hitunglah jumlah 6 suku pertama dari deret $1 + 3 + 9 + \dots$
    \item Tentukan jumlah deret geometri tak hingga dari $16 + 8 + 4 + 2 + \dots$
    \item Sebuah bola dijatuhkan dari ketinggian 10 meter dan memantul kembali dengan ketinggian $3/4$ kali tinggi sebelumnya. Tentukan panjang lintasan bola sampai berhenti.
\end{enumerate}

% ==============================================================
\section{PERSAMAAN DAN FUNGSI KUADRAT}
% ==============================================================
\begin{enumerate}
    \item Tentukan akar-akar dari persamaan $x^2 - 7x + 10 = 0$ dengan faktorisasi.
    \item Gunakan rumus ABC untuk mencari akar dari $2x^2 + 5x - 3 = 0$.
    \item Tentukan nilai Diskriminan ($D$) dari $x^2 - 4x + 4 = 0$ dan tentukan jenis akarnya.
    \item Jika $x_1$ dan $x_2$ adalah akar dari $x^2 + 6x + 5 = 0$, tentukan nilai $x_1 + x_2$ dan $x_1 \cdot x_2$.
    \item Tentukan koordinat titik puncak dari fungsi $f(x) = x^2 - 6x + 8$.
    \item Apakah grafik fungsi $f(x) = -x^2 + 2x + 5$ terbuka ke atas atau ke bawah? Jelaskan.
    \item Tentukan sumbu simetri dari fungsi $f(x) = 2x^2 - 8x + 3$.
    \item Tentukan titik potong terhadap sumbu $y$ dari fungsi $f(x) = 3x^2 + 5x - 12$.
    \item Susunlah persamaan kuadrat yang akar-akarnya adalah 3 dan -5.
    \item Tentukan nilai $k$ agar persamaan $x^2 + kx + 9 = 0$ memiliki akar kembar.
\end{enumerate}

% ==============================================================
\section{TRIGONOMETRI}
% ==============================================================
\begin{enumerate}
    \item Pada segitiga siku-siku $ABC$, jika $\sin A = 3/5$, tentukan nilai $\cos A$ dan $\tan A$.
    \item Hitunglah nilai dari $\sin 30^\circ + \cos 60^\circ - \tan 45^\circ$.
    \item Sebuah tangga panjangnya 6 meter disandarkan pada tembok dan membentuk sudut $60^\circ$ dengan lantai. Berapa tinggi tembok yang dicapai tangga?
    \item Pada segitiga $PQR$, diketahui $p = 8$ cm, $q = 10$ cm, dan sudut $R = 60^\circ$. Tentukan panjang sisi $r$ menggunakan aturan Cosinus.
    \item Gunakan aturan Sinus: Pada segitiga $ABC$, sudut $A = 30^\circ$, sudut $B = 45^\circ$, dan sisi $a = 10$ cm. Tentukan panjang sisi $b$.
    \item Tentukan nilai dari $\sin 120^\circ$.
    \item Sederhanakan $(\sin x)(\csc x)$.
    \item Jika $\tan \theta = 1$ dan $\theta$ berada di kuadran I, tentukan nilai $\theta$.
\end{enumerate}

% ==============================================================
\section{STATISTIKA}
% ==============================================================
\begin{enumerate}
    \item Tentukan Mean, Median, dan Modus dari data: $6, 8, 5, 7, 6, 8, 9, 6, 10$.
    \item Diberikan data frekuensi berikut untuk tinggi badan siswa:
    \begin{center}
    \begin{tabular}{|c|c|}
    \hline
    Tinggi (cm) & Frekuensi \\
    \hline
    150--154 & 4 \\
    155--159 & 10 \\
    160--164 & 6 \\
    165--169 & 8 \\
    170--174 & 2 \\
    \hline
    \end{tabular}
    \end{center}
    Hitunglah nilai \textbf{Mean} dari data kelompok di atas.
    \item Berdasarkan tabel di atas, tentukan letak kelas \textbf{Median}.
    \item Hitunglah \textbf{Modus} dari data pada tabel tersebut.
    \item Jika rata-rata dari data $5, 7, 8, 9, x$ adalah 8, tentukan nilai $x$.
\end{enumerate}

\end{document}